\documentclass[a4paper,10pt]{article}
\newcommand\subdir{/home/koen/LaTeX-setup/}
\input{\subdir/LaTeX-files/imports.tex}
\input{\subdir/LaTeX-files/master-internship.tex}
\usepackage{\subdir/LaTeX-files/aastex}
\usepackage{natbib}
\bibliographystyle{\subdir/LaTeX-files/astroadstitle.bst}

%Set the title, Author and date
\title{Notes Week 10}
\author{Koen Essers}
\tutorial{Master Internship}
\date{\today}

%Set the header style
\header{\tutorialstring}

%Begin the document
\begin{document}
\setuplayout
\section{Stellar Evolution Meeting}
\begin{figure}[ht]
    \centering
    \incplot{w10-stellar-evolution-meeting}
\end{figure}
\newsection{Radiation Corrected Transfer Rate}
As described by \citet{temmink2025}, the MESA setting \lstinline[style=output, breaklines=true]|use_radiation_corrected_transfer_rate| is indeed on by default. From now on, I disable this in the binary inlist.
\newsection{\texorpdfstring{\texttt{evolve.py}}{evolve.py}-script with Rees Single Star TPAGB Models}
In order to 

\newsection{Wind accretion}
% implement the wind accretion
As mentioned last week, we need to take accretion of the secondary into account when we use the wind prescription by \citet{saladino2019}. In this prescription, \(\beta\) describes how much of the wind mass lost from the donor gets accreted by the accretor. Since we already have a nice way to compute this \(\beta \), we can simply use the same methods to find \(\beta \) for the change in mass, and hook this method into the routine \lstinline[style=output, breaklines=true]|other_binary_wind_transfer|. Below, I write the full routine:
\begin{minted}[fontsize=\small]{fortran}
subroutine wind_accretion(binary_id, s_i, ierr)
    integer, intent(in) :: binary_id, s_i ! s_i is index of the wind mass losing star
    integer, intent(out) :: ierr

    type(binary_info), pointer :: b
    type(star_info), pointer :: s

    real(dp) :: q, vw_over_vorb

    call binary_ptr(binary_id, b, ierr)
    if (ierr /= 0) then
        write (*, *) 'failed in binary_ptr'
        return
    end if

    ! --- compute q == m_accretor / m_donor and v
    q = b%m(b%a_i)/b%m(b%d_i)
    vw_over_vorb = compute_vwind(b%s_donor)/compute_vorbit(b%m(1) + b%m(2), b%separation)

    ! --- get beta from Saladino
    beta = compute_beta_Sal(q, vw_over_vorb)/sqrt(1 - pow2(b%eccentricity))
    b%wind_xfer_fraction(s_i) = beta

end subroutine wind_accretion
\end{minted}

It is important to tell MESA to actually compute wind accretion\sidenote{Which I forgot to do initially.}.
Although I linked my custom subroutine to \lstinline[style=output, breaklines=true]|other_binary_wind_transfer| using
\begin{minted}[fontsize=\small,breaklines]{fortran}
b% other_binary_wind_transfer => wind_accretion 
\end{minted}
and set this in the inlist of the binary as well using
\begin{minted}[fontsize=\small,breaklines]{fortran}
use_other_binary_wind_transfer = .true.
\end{minted}
MESA still does not automatically actually use the wind mass transfer routine. For this, one \emph{must} enable the option
\begin{minted}[fontsize=\small,breaklines]{fortran}
do_wind_mass_transfer_1 = .true.
\end{minted}






\newsection{Validating Tides + Wind Prescription}

In order to ensure the custom tides and wind prescriptions are working as expected, I try to thoroughly check every component of the implementation. I start with the computation of the ratio \(k/T\)\sidenote{Which is used in the tidal prescription by \citet{hut1981} and encodes the structure of the star.}. I use the prescription by \citet{preece2022} which uses a power law of the form
\begin{equation*}
    \l(\frac{k}{T}\r)_\text{c} = \l(\frac{R_\text{conv}}{R} \r)^{a} \l(\frac{M_\text{conv}}{M}\r)^{b} \frac{c}{t_\text{conv}}.
\end{equation*}
Clearly, this ratio depends on the mass and radius of the convective envelope, and the so called convective turnover time \(t_\text{conv}\). \citet{preece2022} defines this as
\begin{equation*}
    t_\text{conv} = 0.4311 \l[\frac{3 M_\text{conv} R_\text{conv}^{2} }{L}\r]^{1 / 3}\text{ yr}.
\end{equation*}
I start by looking at \(R_\text{conv}, M_\text{conv}\) and \(t_\text{conv}\) for a reference run.
\begin{figure}[ht]
    \marginfignote{0}{I think it looks good, the radius of the convective envelope is roughly the radius of the star, which is what I would expect. I do however find it a bit weird that the convective envelope seems to become much smaller than the stellar radius during / after RLOF. I want to look into this before I continue.}
    \marginfignote{15}{I therefore investigate some of the profiles of the donor. I label the points of the profiles in the graphs with the open dots. The colors correspond with the lines of the plot below.}
    \centering
    \incplot{w10-checks-2}
\end{figure}

\begin{figure}[ht]
    \marginfignote{0}{The profiles confirm that the convective envelope actually extends a lot less far to the surface as mass-transfer proceeds so the implementation for finding \(R_\text{conv}\) seems to be correct. I am however not sure \emph{why} the superadiabatic region grows in radius during mass-transfer.}
    \centering
    \incplot{w10-check-profile}
\end{figure}
I also show the entropy of the profiles as a function of radius. 
\begin{figure}[ht]
    \marginfignote{0}{Here again, it is quite obvious that the superadiabatic region grows significant in radius.
}
    \centering
    \incplot{w10-check-profile-2}
\end{figure}

\nextpage

Next I check the convective envelope mass:

\begin{figure}[ht]
    \marginfignote{0}{I think it looks good. The mass of the convective region is roughly equal to the envelope mass, and we can even see the core mass grow between pulses. Again, the difference between the star and the convective region increases as RLOF stars, but this time, the effect is very minimal.}
    \centering
    \incplot{w10-checks-3}
\end{figure}

Next, I check \(t_\text{conv}\) and \(f_\text{conv}\), where \(f_\text{conv}\) is a factor related to fast tides.

\begin{figure}[ht]
    \marginfignote{0}{The convective turnover timescale seems to have realistic values. The sharp drop towards the end of the evolution is caused by the decreasing mass with similar radii and luminosities.}
    \marginfignote{12}{The fast tides factor is 1 everywhere, which is the expected result for giant stars where the tidal period is much longer than the convective turnover time.}
    \centering
    \incplot{w10-checks-4}
\end{figure}

Next, the tides depend on the gyration radius \(r_{g}\) and the moment of inertia \(I\). Since these are closely related, I plot them together below.

\begin{figure}[ht]
    \marginfignote{1}{It makes sense that the moment of inertia follows the curve of \(R^{2}\) really closely, since the mass barely changes. This however no longer holds once will mass loss becomes important, and especially, when RLOF starts. A quick estimation of the moment of inertia of a homogenous sphere with \(M= 2M_\odot \) and \(R = 300R_\odot \) gives \(I \approx 7\cdot 10^{59}\) g cm\(^{2}\), which is the same order of magnitude.}
    \marginfignote{16}{The gyration radius encodes how centered the mass is concentrated in the star. For a giant star, it makes sense that \(r_{g}^{2} \approx 0.1\).}
    \centering
    \incplot{w10-checks-10}
\end{figure}


\begin{figure}[ht]
    \centering
    \incplot{w10-checks-1}
\end{figure}





\newsection{Grid}
% run a grid with same initial parameters, but with wind and tides

\newsection{Analytical expressions of grid results}

The goal of this section is to understand the results from the following plot using analytical descriptions of orbital dynamics and mass transfer:
\begin{figure}[ht]
    \centering
    \incplot{w8-grid-2}
\end{figure}

Let's start with the general expression for angular momentum of a two body problem of point masses in the centre of mass coordinate system:
\begin{align*}
    \vec{J} &= \mu \vec{r}\times  \vec{v} = \mu r^{2} \dot{\phi }  \; \hat{r} \times  \hat{\phi }.\\
    \intertext{Here, \(\mu\) is the reduced mass defined as}
    \mu &= \frac{M_{1}M_{2}}{M_{1} + M_{2}}.\\
    \intertext{The magnitude of the orbital angular momentum is then}
    J &= \mu r^{2} \dot{\phi } = \text{const}.\\
    \intertext{Energy is a conserved quantity as well, which we can express as}
    E &= K+ V = \frac{1}{2}\mu \dot{ \vec{r}}^{2} - G \frac{M_{1}M_{2}}{r}\\
      &= \frac{1}{2}\mu \dot{r}^{2} + \frac{1}{2} \mu r^{2} \dot{ \phi }^{2} - G \frac{M_{1}M_{2}}{r}.\\
      \intertext{By substituting the equation for the angular momentum, we find that}
    E &= \frac{1}{2}\mu \dot{r}^{2} + \frac{J^{2} }{2 \mu r^{2}}  - G \frac{M_{1}M_{2}}{r}.\\
    \intertext{Combining this with the standard result for the specific energy of an elliptical orbit,}
    \varepsilon &= - G \frac{M_{1} + M_{2} }{2a} \implies E = \varepsilon \mu = - G \frac{M_{1}M_{2}  }{2a},\\
    \intertext{we can solve for \(J^{2}\):}
    - G \frac{M_{1}M_{2}  }{2a} &= \frac{1}{2}\mu \dot{r}^{2} + \frac{J^{2} }{2 \mu r^{2}},\\
    \frac{J^{2} }{2 \mu r^{2}} &= GM_{1}M_{2} \l[\frac{1}{r} - \frac{1}{2a}\r] - \frac{1}{2} \mu \dot{r}^{2}.\\
    J^{2} &= 2 G r\frac{M_{1}^{2} M_{2}^{2}}{M_{1} + M_{2} }\l[1 - \frac{r}{2a}\r] - \mu^{2} r^{2}  \dot{r}^{2}.\\
    \intertext{Using that total energy is conserved throughout the orbit, we can solve for \(J^{2}\) at one point in the orbit, and the result will hold throughout the orbit. At the periapsis or apoasis, \(\dot{r}=0\), which reduces the equation to}
    J^{2} &= 2 G r_\text{p}\frac{M_{1}^{2} M_{2}^{2}}{M_{1} + M_{2} }\l[1 - \frac{r_\text{p}}{2a}\r].\\
    \intertext{Using \(r_\text{p} = a(1 - e)\), we find that}
    J^{2} &= 2 G \frac{M_{1}^{2} M_{2}^{2}}{M_{1} + M_{2} }\l[\frac{1 + e}{2}\r] a(1 - e).\\
    J^{2} &= G \frac{M_{1}^{2} M_{2}^{2}}{M_{1} + M_{2} }a(1 -e^{2}).
\end{align*}
These are all terms we \emph{know} so we can use this definition of \(J\) as a conserved quantity and explore what happens when we vary masses due to mass transfer. It is important to note that this definition assumes the angular momentum of the binary is strictly the orbital angular momentum of the stars. For our model without tides (and thus no stellar rotation), this is perfectly fine but the upcoming results will not strictly be true for rotating stars.

Let's start by taking the time derivative of \(J^{2}\):
\begin{align*}
    \dv{}{t}(J^{2}) &= 2 J \dv{J}{t} = \dv{}{t}\l[G \frac{M_{1}^{2} M_{2}^{2}}{M_{1} + M_{2} }a(1 -e^{2})\r].\\
    \intertext{It is actually easier to compute \(\dv{}{t}\l(J^{2}\r) / J^{2}\), since a lot of the terms cancel.}
    \frac{2}{J}\dv{J}{t} &= \frac{1}{J^{2}}\dv{}{t}\l[G \frac{M_{1}^{2} M_{2}^{2}}{M_{1} + M_{2} }a(1 -e^{2})\r],\\
    \frac{2}{J}\dv{J}{t} &= \frac{1}{J^{2}}\l\{\dv{}{t}\l[M_{1}^{2}\r] \frac{J^{2}}{M_{1}^{2}}  + \dv{}{t}\l[M_{2}^{2}\r] \frac{J^{2}}{M_{2}^{2}} + \dv{}{t}\l[(M_{1} + M_{2})^{-1} \r]  \cdot  \r.\\
                         & \phantom{=\frac{1}{J^{2}}\;\;\;\;\;}J^{2}\l(M_{1} + M_{2}\r) + \l.\dv{}{t}\l[a\r] \frac{J^{2}}{a} + \dv{}{t}\l[1- e^{2}\r] \frac{J^{2}}{1-e^{2} } \r\},\\
    \frac{2}{J}\dv{J}{t} &= \frac{2 \dot{M}_\text{1}}{M_{1}} + \frac{2 \dot{M}_\text{2}}{M_{2}} - \frac{\dot{M}_{1} + \dot{M}_{2}}{M_{1} + M_{2}} + \frac{\dot{a}}{a} -\frac{2 e\dot{e}}{1-e^{2}}.
\end{align*}
In the case of conservative mass transfer in a circular orbit and no angular momentum loss due to wind mass loss, \(\dot{J}  = 0\), \(e =0\) and \(\dot{M}_{1} = - \dot{M}_{2}\). This allows us to solve for the change in separation, which is closely related to the period.
\begin{align*}
    \frac{\dot{a}}{a} &= \frac{2 \dot{M}_\text{d}}{M_\text{a}} - \frac{2 \dot{M}_\text{d}}{M_\text{d}} = 2 \frac{\dot{M}_\text{d} }{M_\text{d}}\l(\frac{M_\text{d}}{M_\text{a}} - 1\r).\\
    \intertext{Since mass is conserved, \(M_{a} = M_\text{tot} - M_\text{d} \) is always true and \(M_\text{tot}\) is constant. This means we can write}
    \frac{\dot{a}}{a} &= \frac{2\dot{M}_\text{d} }{M_\text{d}}\l(\frac{M_\text{d}}{M_\text{tot} - M_\text{d}} - 1\r),\\
    &= 2 \frac{\dot{M}_\text{d} }{M_\text{d}}\l(\frac{2M_\text{d} - M_\text{tot}}{M_\text{tot} - M_\text{d}}\r),\\
    \frac{\dd{a}}{a} &= \frac{2}{M_\text{d}} \l(\frac{2M_\text{d} - M_\text{tot}}{M_\text{tot} - M_\text{d}} \r) \dd{M}_\text{d},\\
    \int_{a_\text{i}}^{a_\text{f}} \frac{\dd{a}}{a} &= \int_{M_\text{d,i}}^{M_\text{d,f}} \frac{2}{M_\text{d}} \l(\frac{2M_\text{d} - M_\text{tot}}{M_\text{tot} - M_\text{d}} \r)  \,\dd{M_\text{d}},\\
    \ln \l(\frac{a_\text{f}}{a_\text{i}}\r)  &= 2\int_{M_\text{d,i}}^{M_\text{d,f}} \frac{2M_\text{d} - M_\text{tot}}{M_\text{d} \l(M_\text{tot} - M_\text{d}\r) } \,\dd{M_\text{d}},\\
    \intertext{Splitting the fraction yields}
    \ln \l(\frac{a_\text{f}}{a_\text{i}}\r)  &= 2\int_{M_\text{d,i}}^{M_\text{d,f}}-\frac{1}{M_\text{d}} + \frac{1}{M_\text{tot} - M_\text{d}} \,\dd{M_\text{d}},\\
    \ln \l(\frac{a_\text{f}}{a_\text{i}}\r)  &= 2 \ln \l(\frac{M_\text{d,i}}{M_\text{d,f}}\r) + 2\ln \l(\frac{M_\text{tot} - M_\text{d,i}}{M_\text{tot}-M_\text{d,f}}\r)\\
    \frac{a_\text{f}}{a_\text{i}} &= \l(\frac{M_\text{d,i}\l(M_\text{tot} - M_\text{d,i}\r) }{M_\text{d,f}\l(M_\text{tot} - M_\text{d,f}\r) }\r)^{2}.
\end{align*}
Cool, Let's show this relation in a couple of different ways:
\begin{figure}[ht]
    \centering
    \incplot{w10-analytical1}
\end{figure}

\begin{figure}[ht]
    \centering
    \incplot{w10-analytical2}
\end{figure}



\bibliography{master.bib}
\end{document}
